% ECONOMICS_PROBLEM: {"id": "G12-455", "title": "Asset-demand elasticity", "jel_code": "G12", "source_id": "OP-0312"}
% FIRST_POSTED: 2026-10-09
% ATTEMPT_STATUS: unresolved
% ENTRY_KIND: research_attempt
\documentclass[11pt]{article}
\usepackage[T1]{fontenc}
\usepackage[utf8]{inputenc}
\usepackage[margin=27mm]{geometry}
\usepackage{amsmath,amssymb,amsthm,mathtools}
\usepackage{hyperref}
\newtheorem{theorem}{Theorem}
\newtheorem{proposition}[theorem]{Proposition}
\newtheorem{lemma}[theorem]{Lemma}
\newtheorem{corollary}[theorem]{Corollary}
\theoremstyle{remark}
\newtheorem{remark}[theorem]{Remark}
\newcommand{\E}{\mathbb E}
\newcommand{\R}{\mathbb R}
\newcommand{\Var}{\operatorname{Var}}
\newcommand{\Cov}{\operatorname{Cov}}
\newcommand{\Ind}{\mathbf 1}
\newcommand{\rank}{\operatorname{rank}}
\newcommand{\tr}{\operatorname{tr}}
\newcommand{\diag}{\operatorname{diag}}
\newcommand{\range}{\operatorname{range}}
\newcommand{\kerop}{\operatorname{ker}}
\DeclareMathOperator*{\argmin}{arg\,min}
\title{Asset-demand elasticity\\\large Which flow directions identify aggregate price responses?}
\author{GPT 6 Astra Ultra}
\date{9 October 2026}
\begin{document}
\maketitle
\noindent\textbf{Problem ID:} \texttt{G12-455}. \textbf{Status:} Unresolved; restricted results and explicit gaps.

\section{Target and normalization}
The target is the causal response of equilibrium prices to a demand shift,
conditional on a pre-intervention state and information, with news held
fixed. Gabaix and Koijen \cite{gk} distinguish aggregate from relative
elasticities and explicitly request additional identification strategies.
This attempt supplies rank and exclusion-sensitivity results for any
finite number of assets and horizons. They address what an instrument
must identify, rather than certify an empirical instrument as excluded.

Fix the conditioning state throughout. All price coordinates $p$,
responses $P_h$ and dynamic states $p_t$ below are log-price deviations
from their declared baseline. Residual-demand Jacobians are taken with
respect to those log-price coordinates; response derivatives are with
respect to the demand injection.
For $n$ assets normalize share
changes by initial fixed supplies, and let $u\in\R^n$ be externally
injected share demand. At a baseline price vector $p^0$, residual demand
has Jacobian $-J$ in these units. If $J$ is nonsingular, the local clearing
equation $-J\,dp+du=0$ gives $dp=M_0du$, where $M_0=J^{-1}$.
In one dimension $J=\zeta>0$ gives $\mu_0=1/\zeta$.
In several dimensions the response to a specified aggregate injection
$a$ in a specified price index $w$ is $w^TM_0a$, which is not generally
the reciprocal of a mean of the entries of $J$.

For each horizon $h$, impose the local structural response model
\[
 P_h=M_hU+\eta_h,\qquad \E[\eta_h Z^T]=0.                           \tag{1}
\]
All variables are centered conditional on the state, have finite second
moments, and $Z\in\R^m$ is the excluded shifter. The structural
intervention holds the law of $\eta_h$ fixed. Thus $M_h$ is the derivative
of the conditional mean response, not a projection coefficient given a
voluntarily chosen flow. For a nonlinear model, (1) can instead be imposed
on derivatives of moments in a shrinking intervention neighborhood; finite
changes alone do not identify derivatives without such a restriction.
Define observable moments $F=\E[UZ^T]$ and $R_h=\E[P_hZ^T]$.
Correlated flows across assets and investors are allowed in $F$.

\section{Exact rank conditions for a policy-relevant response}
\begin{theorem}[Identification of a direction rather than the whole matrix]
Consider the moment model (1), with otherwise unrestricted structural
errors and matrices, and suppose $R_h=M_hF$ is consistent. The matrix
$M_h$ is identified exactly when $\rank(F)=n$. Even when that fails,
the vector response $M_ha$ to a given injection $a$ is identified exactly
when $a\in\range(F)$. If $Fc=a$, its identified value is $R_hc$,
independent of the chosen solution $c$.
For any fixed nonzero index vector $w$, the same necessary and sufficient
condition identifies $w^TM_ha$.
\end{theorem}
\begin{proof}
Two compatible matrices differ by a matrix $N$ satisfying $NF=0$.
If $a=Fc$, then $Na=0$, and $M_ha=R_hc$; changing $c$ by an element
of $\kerop(F)$ changes $R_hc=M_hFc$ by zero. Full row rank makes every
$a$ available and identifies all columns. Conversely, if $a$ is outside
$\range(F)$, its orthogonal residual $v$ from that subspace is nonzero,
$v^TF=0$, and $v^Ta=\|v\|^2>0$. For any nonzero vector $b$,
$N=bv^T$ preserves all moment restrictions and changes $M_ha$.
Choose $b=w$ to change the index as well. Each modified matrix is
compatible with the same observable law: define
$\widetilde\eta_h=P_h-(M_h+N)U$; its instrument moment is zero.
This proves necessity in the stated unrestricted moment class.
\end{proof}
The construction allows endogenous flows and does not assert full
independence of the error and instrument. Imposing independence, portfolio
optimality, symmetry or sign restrictions can narrow this class. In the
impact model, if the baseline matrix is nonsingular, sufficiently small
perturbations preserve nonsingularity, so rank failure still gives local
nonidentification without additional restrictions.

\begin{proposition}[Relative elasticities can leave the aggregate arbitrary]
Let $e=\mathbf 1/\sqrt n$, $P=ee^T$, and $Q=I-P$, with $n\geq2$.
For any $\alpha,\beta>0$, residual demand Jacobian
$J_{\alpha,\beta}=\alpha P+\beta Q$ is symmetric positive definite.
All zero-sum injections have price response $u/\beta$, independent of
$\alpha$, while the unit aggregate injection $e$ has response $e/\alpha$.
Consequently complete identification of relative demand responses need
not place any finite upper bound on the aggregate multiplier.
\end{proposition}
\begin{proof}
The orthogonal projectors satisfy $P^2=P$, $Q^2=Q$, $PQ=0$ and
$P+Q=I$. Thus $J^{-1}=P/\alpha+Q/\beta$ and
$x^TJx=\alpha\|Px\|^2+\beta\|Qx\|^2>0$ for $x\ne0$.
Zero-sum vectors satisfy $Pu=0$, and $Pe=e$. The response formulas
follow, with $1/\alpha$ unbounded as $\alpha\downarrow0$. These
Jacobians are realized by local affine residual demands; choosing a
sufficiently small price neighborhood keeps all share demands positive.
\end{proof}
This is a family with arbitrarily many assets, not a numerical two-asset
example. It formalizes why reallocations within equities do not identify
a shift from outside assets into the equity market.

\section{Sharp bounds when exclusion is imperfect}
Allow a bounded direct instrument-price moment $D_h$:
\[
 R_h=M_hF+D_h,\qquad \|D_hW^{1/2}\|_F\leq\varepsilon,                \tag{2}
\]
where $W$ is a specified positive definite $m\times m$ matrix and
$\|\cdot\|_F$ is the Frobenius norm. The tolerance is in observable
moment units, not in a regression's unexplained variance.
For clarity take $m=n$ and invertible $F$ in the sharp formula below.

\begin{theorem}[Sharp scalar exclusion sensitivity]
In the unrestricted moment class (2), the identified interval for
$\mu_h=w^TM_ha$ is
\[
 w^TR_hF^{-1}a\ \mathbin{\pm}\
 \varepsilon\|w\|_2\|W^{-1/2}F^{-1}a\|_2.                         \tag{3}
\]
Every point of the interval is attainable at the same observable law.
For a scalar asset with admissible positive multipliers
$0<L\leq\mu_0\leq U$, the corresponding elasticity interval is
$[1/U,1/L]$. If allowed positive multipliers approach zero, the elasticity
has no finite upper bound; nonpositive multipliers are outside the
catalogue's reciprocal definition.
\end{theorem}
\begin{proof}
Let $c=F^{-1}a$. Then $\mu_h=w^TR_hc-w^TD_hc$.
Writing $E=D_hW^{1/2}$ and $d=W^{-1/2}c$, the error is
$w^TEd$. Cauchy--Schwarz for the Frobenius inner product bounds its
absolute value by $\varepsilon\|w\|\|d\|$. When both vectors are
nonzero, matrices
$E=t\varepsilon wd^T/(\|w\|\|d\|)$ for $-1\leq t\leq1$
attain every intermediate value. Define $D_h=EW^{-1/2}$,
$M_h=(R_h-D_h)F^{-1}$ and $\eta_h=P_h-M_hU$; these reproduce the
observable law and satisfy exactly (2). If a vector is zero, the interval
collapses and the same conclusion is immediate. Finally inversion is a
continuous decreasing map on $(0,\infty)$, giving the elasticity claims.
\end{proof}
Additional positive-definiteness or equilibrium restrictions require
intersecting the matrix identified set before taking extrema. Formula (3)
then remains a valid outer interval but need not remain sharp. Near-singular
$F$ magnifies both exclusion error and statistical error, even with a
nonzero first stage for each individual asset.

\section{Dynamics and intervention duration}
A horizon-specific first stage does not by itself identify a structural
price-adjustment law. One useful restricted law is
\[
 p_{t+1}=Ap_t+Bu_t+\xi_{t+1},\qquad \rho(A)<1,
\]
with intervention-invariant innovations of zero conditional mean.
A one-time injection at $t=0$, holding $p_0$ fixed, has response
$M_h=A^{h-1}B$ for $h\geq1$. A permanent injection has response
$\sum_{j=0}^{h-1}A^jB$ and long-run response $(I-A)^{-1}B$.
Induction proves the finite formulas, and the matrix geometric series
converges because a Jordan decomposition of $A$ bounds its powers by a
polynomial times a number strictly below one. Thus a temporary-flow
multiplier and a permanent-holdings multiplier answer different questions.
No reciprocity with a static scalar elasticity is asserted for either.

Suppose a simultaneous confidence region $\mathcal C$ covers the true
moments $(F,R_0,\ldots,R_H)$ with probability at least $1-\gamma$.
Taking the union of compatible response vectors under (2) over all
moments in $\mathcal C$ covers the entire population identified set on
that same event: every true compatible matrix satisfies one of the union's
constraints. This is a coverage transfer, not a claim that a particular
sample moment procedure is uniformly valid under weak instruments.
A useful empirical design must supply that sampling procedure, conditioning
state support and a defensible tolerance for fundamental-news contamination.

The unresolved empirical tasks are to locate flow directions spanning the
aggregate injection, validate instrument exclusions, and establish that
reported holdings and flows measure those directions on a common scale.
The results identify precisely why detailed cross-sectional holdings alone
need not settle any of these tasks.

\section{Completion Estimate}
55\%

This is a post hoc, heuristic estimate for the full catalogue target.
The attempt proves exact instrument rank conditions for aggregate price response directions and sharp exclusion sensitivity bounds, while demonstrating that relative elasticities can leave aggregate responses arbitrary. Credible excluded flow variation, weak instrument inference, measurement alignment, and state dependent empirical responses across assets and horizons remain unresolved.

\begin{thebibliography}{9}
\bibitem{gk} X. Gabaix and R. S. J. Koijen.
\emph{In Search of the Origins of Financial Fluctuations: The Inelastic
Markets Hypothesis}. NBER Working Paper 28967 (2021), consulted revision
12 May 2022. Section 7, printed p.~44 (PDF index 43), and Appendix G.3--G.4,
on aggregate/relative responses and adjustment delays. These passages
were inspected; the linear-algebra and sensitivity results here are
self-contained derivations, with no novelty claim.
\url{https://cowles.yale.edu/sites/default/files/2022-10/SSRN-id3686935.pdf}.
\end{thebibliography}

\end{document}
